\documentclass[a4paper,12pt]{article}

\usepackage[T2A]{fontenc}
\usepackage[utf8]{inputenc}
\usepackage[russian]{babel}
\usepackage{amsmath,amssymb,mathtools}
\usepackage{PTSans}
\renewcommand{\familydefault}{\sfdefault}
\usepackage{icomma}
\usepackage[
  a4paper,
  left=18mm,
  right=18mm,
  top=18mm,
  bottom=18mm,
  headheight=25pt,
  headsep=6mm,
  footskip=11mm
]{geometry}
\usepackage{microtype}
\usepackage{tikz}
\usetikzlibrary{calc,arrows.meta,positioning,shapes.geometric}
\usepackage{tabularx,booktabs}
\usepackage{enumitem}
\usepackage{parskip}
\usepackage{fancyhdr}
\usepackage{setspace}
\usepackage{xcolor}

\definecolor{accent}{HTML}{008080}
\definecolor{darkaccent}{HTML}{004D4D}
\definecolor{textgray}{HTML}{333333}
\definecolor{softgray}{HTML}{6B7280}

\onehalfspacing
\setlength{\parindent}{0pt}
\setlength{\emergencystretch}{2em}
\raggedbottom

\pagestyle{fancy}
\fancyhf{}
\lhead{\small\color{softgray}Володя Хачатурян --- ОГЭ}
\rhead{\small\color{softgray}Средняя скорость}
\cfoot{\small\color{softgray}\thepage}
\renewcommand{\headrulewidth}{0.3pt}
\renewcommand{\headrule}{%
  \hbox to\headwidth{%
    \color{accent}\leaders\hrule height \headrulewidth\hfill%
  }%
}

\newcommand{\sectionline}{%
  \vspace{0.15em}%
  \noindent\textcolor{accent}{\rule{\linewidth}{0.6pt}}%
  \vspace{0.3em}%
}

\newcommand{\checkitem}[1]{\item[$\square$] #1}

\tikzset{
  flowstart/.style={
    ellipse,
    draw=darkaccent,
    line width=0.9pt,
    align=center,
    minimum width=40mm,
    minimum height=10mm,
    font=\small
  },
  flowproc/.style={
    rectangle,
    rounded corners=2mm,
    draw=accent,
    line width=0.9pt,
    align=center,
    text width=98mm,
    minimum height=13mm,
    font=\small,
    inner sep=3mm
  },
  flowdec/.style={
    diamond,
    aspect=2.35,
    draw=darkaccent,
    line width=0.9pt,
    align=center,
    text width=58mm,
    font=\small,
    inner sep=2mm
  },
  flowline/.style={
    -{Latex[length=3mm,width=2mm]},
    line width=0.9pt,
    draw=textgray
  }
}

\begin{document}

% ============================================================
% ТИТУЛЬНЫЙ ЛИСТ
% ============================================================

\begin{titlepage}
\thispagestyle{empty}
\centering

\vspace*{18mm}

{\color{darkaccent}\Large\bfseries Чек-лист\par}

\vspace{5mm}

{\huge\bfseries Средняя скорость и базовая кинематика\par}

\vspace{7mm}

{\Large Володя Хачатурян\par}

\vspace{2mm}

{\large ОГЭ по физике\par}

\vspace{14mm}

\textcolor{accent}{\rule{0.72\linewidth}{0.9pt}}\par

\vspace{8mm}

{\large Путь Б: 27.09--28.09.2026\par}

\vspace{4mm}

{\large \today\par}

\vfill

{\normalsize
Лёвин Артём Александрович эксклюзивно для Володи Хачатуряна
\par}

\vspace{15mm}

\begin{tikzpicture}[remember picture,overlay]
  \draw[
    accent,
    line width=1.1pt
  ]
    ($(current page.south west)+(16mm,13mm)$)
      .. controls
      ($(current page.south west)+(62mm,24mm)$)
      and
      ($(current page.south east)+(-62mm,6mm)$)
      ..
      ($(current page.south east)+(-16mm,15mm)$);
\end{tikzpicture}

\end{titlepage}

% ============================================================
% БАЗОВОЕ ПОВТОРЕНИЕ
% ============================================================

\section*{1. Базовая кинематика: что нужно помнить}
\sectionline

\subsection*{1.1. Координата, перемещение, скорость, ускорение}

Для равноускоренного движения вдоль оси $Ox$:

\[
x=x_0+v_{0x}t+\frac{a_xt^2}{2},
\qquad
v_x=v_{0x}+a_xt.
\]

Изменение координаты и проекция перемещения:

\[
\Delta x=x-x_0.
\]

Если движение происходит вдоль одной прямой и тело на рассматриваемом промежутке не меняет направление, то

\[
S=|\Delta x|.
\]

Если тело меняло направление, путь следует находить как сумму длин всех пройденных участков.

Ускорение:

\[
a_x=\frac{\Delta v_x}{\Delta t}
=\frac{v_x-v_{0x}}{t-t_0}.
\]

Если отсчёт времени начинается с $t_0=0$, то

\[
a_x=\frac{v_x-v_{0x}}{t}.
\]

Для равноускоренного движения полезны три формы перемещения:

\[
\boxed{\Delta x=v_{0x}t+\frac{a_xt^2}{2}},
\]

\[
\boxed{\Delta x=\frac{v_{0x}+v_x}{2}\,t},
\]

\[
\boxed{\Delta x=\frac{v_x^2-v_{0x}^2}{2a_x}},
\qquad a_x\ne0.
\]

Последняя формула удобна, когда в условии нет времени.

\subsection*{1.2. Путь за несколько секунд и за одну $n$-ю секунду}

Пусть $\tau=1\text{ с}$. При постоянном ускорении путь за первые $n$ секунд, если направление движения не меняется:

\[
S_{1\ldots n}=v_0n\tau+\frac{a(n\tau)^2}{2}.
\]

Путь, пройденный только за $n$-ю секунду, удобно получить как разность:

\[
S_n=S(n\tau)-S((n-1)\tau).
\]

Отсюда

\[
\boxed{S_n=v_0\tau+\frac{a\tau^2}{2}(2n-1)},
\qquad \tau=1\text{ с}.
\]

Так сохраняется размерностная корректность формулы: оба слагаемых имеют размерность длины.

\subsection*{1.3. Короткое повторение геометрии}

\begin{tabularx}{\linewidth}{@{}>{\bfseries}l X@{}}
\toprule
Прямоугольный треугольник
&
$c^2=a^2+b^2$, где $c$ --- гипотенуза.
\\
Длина окружности
&
$L=2\pi R=\pi D$, где $D=2R$.
\\
Половина окружности
&
$L_{1/2}=\pi R$.
\\
Четверть окружности
&
$L_{1/4}=\dfrac{\pi R}{2}$.
\\
\bottomrule
\end{tabularx}

% ============================================================
% СРЕДНЯЯ СКОРОСТЬ
% ============================================================

\section*{2. Средняя скорость: основная идея}
\sectionline

В задачах этого чек-листа под средней скоростью понимается \textbf{средняя путевая скорость}. Она отвечает на вопрос: с какой одной постоянной скоростью нужно было бы пройти тот же полный путь за то же полное время?

Главная формула:

\[
\boxed{
v_{\text{ср}}=
\frac{S_{\Sigma}}{t_{\Sigma}}=
\frac{S_1+S_2+\dots+S_n}{t_1+t_2+\dots+t_n}
}
\]

Для каждого участка равномерного движения:

\[
S_i=v_it_i,
\qquad
t_i=\frac{S_i}{v_i},
\qquad v_i>0.
\]

\textbf{Главный принцип:} сначала восстановите полный путь и полное время, затем применяйте формулу средней скорости.

\subsection*{2.1. Известны время и скорость на каждом участке}

\textbf{Пример.}
Первые $5$ ч тело двигалось со скоростью $60$ км/ч, следующие $4$ ч --- со скоростью $15$ км/ч.

\[
S_1=60\cdot5=300\text{ км},
\qquad
S_2=15\cdot4=60\text{ км}.
\]

\[
v_{\text{ср}}=
\frac{300+60}{5+4}
=40\text{ км/ч}.
\]

\subsection*{2.2. Равные части пути}

Пусть первую половину пути тело прошло со скоростью $v_1$, вторую половину --- со скоростью $v_2$. Тогда

\[
S_1=S_2=\frac{S}{2}.
\]

Время на каждом участке:

\[
t_1=\frac{S}{2v_1},
\qquad
t_2=\frac{S}{2v_2}.
\]

Поэтому

\[
v_{\text{ср}}=
\frac{S}{\dfrac{S}{2v_1}+\dfrac{S}{2v_2}}
=\boxed{\frac{2v_1v_2}{v_1+v_2}}.
\]

\textbf{Пример.}
Половину пути человек шёл со скоростью $5$ км/ч, вторую половину --- со скоростью $10$ км/ч.

\[
v_{\text{ср}}=
\frac{2\cdot5\cdot10}{5+10}
=\frac{20}{3}\text{ км/ч}
\approx6{,}67\text{ км/ч}.
\]

Здесь нельзя брать обычное среднее арифметическое скоростей: на более медленном участке человек проводит больше времени.

\subsection*{2.3. Равные части времени}

Если первую половину времени тело двигалось со скоростью $v_1$, вторую половину --- со скоростью $v_2$, то

\[
t_1=t_2=\frac{T}{2}.
\]

Тогда

\[
S_1=v_1\frac{T}{2},
\qquad
S_2=v_2\frac{T}{2},
\]

и

\[
v_{\text{ср}}=
\frac{v_1\dfrac{T}{2}+v_2\dfrac{T}{2}}{T}
=\boxed{\frac{v_1+v_2}{2}}.
\]

\textbf{Пример.}
Половину времени автомобиль двигался со скоростью $20$ км/ч, вторую половину --- со скоростью $80$ км/ч.

\[
v_{\text{ср}}=
\frac{20+80}{2}
=50\text{ км/ч}.
\]

\subsection*{2.4. Доли и проценты пути}

Процент пути переводится в долю полного пути:

\[
20\%\,S=\frac{20}{100}S=\frac{S}{5},
\]

\[
50\%\,S=\frac{S}{2},
\qquad
30\%\,S=\frac{3S}{10}.
\]

Если доли пути равны $p_1,p_2,\dots,p_n$ и

\[
p_1+p_2+\dots+p_n=1,
\]

то

\[
S_i=p_iS,
\qquad
t_i=\frac{p_iS}{v_i}.
\]

Следовательно,

\[
\boxed{
v_{\text{ср}}=
\frac{1}{\dfrac{p_1}{v_1}+\dfrac{p_2}{v_2}+\dots+\dfrac{p_n}{v_n}}
}
\]

при $v_1,v_2,\dots,v_n>0$.

\textbf{Пример.}
$20\%$ пути тело двигалось со скоростью $10$ м/с, $50\%$ --- со скоростью $12$ м/с, оставшиеся $30\%$ --- со скоростью $15$ м/с.

\[
v_{\text{ср}}=
\frac{1}{\dfrac{0{,}2}{10}+\dfrac{0{,}5}{12}+\dfrac{0{,}3}{15}}.
\]

\[
\frac{0{,}2}{10}=\frac{1}{50},
\qquad
\frac{0{,}5}{12}=\frac{1}{24},
\qquad
\frac{0{,}3}{15}=\frac{1}{50}.
\]

\[
v_{\text{ср}}=
\frac{1}{\dfrac{1}{50}+\dfrac{1}{24}+\dfrac{1}{50}}
=\frac{600}{49}\text{ м/с}
\approx12{,}24\text{ м/с}.
\]

% ============================================================
% АЛГЕБРА
% ============================================================

\section*{3. Алгебраический минимум для этих задач}
\sectionline

\subsection*{3.1. Деление дроби на число}

Если дробь дополнительно делится на ненулевое число, это число становится множителем знаменателя:

\[
\frac{A}{B}:c=\frac{A}{Bc},
\qquad
B\ne0,
\quad
c\ne0.
\]

Например,

\[
\frac{S}{2}:5=\frac{S}{10},
\qquad
\frac{S}{3}:10=\frac{S}{30}.
\]

Именно это действие возникает при вычислении времени:

\[
t=\frac{S}{v}.
\]

\subsection*{3.2. Сложение дробей}

Для двух дробей можно сразу использовать формулу

\[
\frac{A}{B}+\frac{C}{D}
=\frac{AD+BC}{BD},
\qquad
B\ne0,
\quad
D\ne0.
\]

Например,

\[
\frac{1}{5}+\frac{2}{3}
=\frac{1\cdot3+2\cdot5}{5\cdot3}
=\frac{13}{15}.
\]

В задачах на среднюю скорость часто удобно сначала привести дроби к общему знаменателю:

\[
\frac{S}{10}+\frac{S}{20}
=\frac{2S}{20}+\frac{S}{20}
=\frac{3S}{20}.
\]

Если дробей три и больше, общий знаменатель обычно делает вычисления понятнее и надёжнее.

\subsection*{3.3. Деление на дробь}

При делении на дробь умножаем на обратную дробь:

\[
\frac{A}{B}:\frac{C}{D}
=\frac{A}{B}\cdot\frac{D}{C}
=\frac{AD}{BC},
\]

где

\[
B\ne0,
\qquad
C\ne0,
\qquad
D\ne0.
\]

Пример:

\[
\frac{1}{3}:\frac{2}{5}
=\frac{1}{3}\cdot\frac{5}{2}
=\frac{5}{6}.
\]

Для выражения, типичного в задачах на среднюю скорость,

\[
\frac{S}{\dfrac{3S}{20}}
=\frac{S}{1}\cdot\frac{20}{3S}
=\frac{20}{3},
\qquad S\ne0.
\]

В физической задаче о пройденном пути обычно $S>0$, поэтому сокращение допустимо.

\subsection*{3.4. Сокращение алгебраических дробей}

Сокращать можно только общие \textbf{множители}.

\[
\frac{5S}{10S}=\frac12,
\qquad S\ne0.
\]

В выражении

\[
\frac{S+5}{S}
\]

сокращение на $S$ невозможно: числитель представляет собой сумму.

Если же

\[
\frac{S(S+5)}{S},
\qquad S\ne0,
\]

то общий множитель $S$ сокращается:

\[
\frac{S(S+5)}{S}=S+5.
\]

% ============================================================
% ТИПИЧНЫЕ ОШИБКИ
% ============================================================

\section*{4. Типичные ошибки и быстрые проверки}
\sectionline

\begin{enumerate}[leftmargin=8mm,itemsep=4pt]

\item
\textbf{Брать среднее арифметическое любых двух скоростей.}
Формула
\[
\frac{v_1+v_2}{2}
\]
работает для двух равных промежутков времени. Для двух равных половин пути используется
\[
\frac{2v_1v_2}{v_1+v_2}.
\]

\item
\textbf{Путать путь и перемещение.}
При движении с разворотом путь складывается по участкам и всегда неотрицателен.

\item
\textbf{Не переводить условие на математический язык.}
Например,
\[
\text{половина пути}\Rightarrow\frac{S}{2},
\qquad
\text{две трети времени}\Rightarrow\frac{2T}{3},
\qquad
20\%\text{ пути}\Rightarrow0{,}2S.
\]

\item
\textbf{Ошибаться при делении дроби на число.}
Правильно:
\[
\frac{S}{2}:5=\frac{S}{10}.
\]

\item
\textbf{Складывать дроби без общего знаменателя.}
Правильно:
\[
\frac{S}{10}+\frac{S}{20}=\frac{3S}{20}.
\]

\item
\textbf{Складывать скорости в знаменателе формулы средней скорости.}
Главная формула всегда имеет структуру
\[
\boxed{v_{\text{ср}}=\frac{\text{весь путь}}{\text{всё время}}}.
\]

\item
\textbf{Сокращать слагаемые.}
Сокращение допустимо только для общих множителей при ненулевом сокращаемом выражении.

\item
\textbf{Смешивать единицы измерения.}
Перед сложением все пути должны быть выражены в одних единицах, все времена --- также в одних единицах.

\end{enumerate}

\textbf{Быстрая проверка ответа.}
Если все скорости на участках положительны, средняя скорость должна находиться между минимальной и максимальной скоростями этих участков.

% ============================================================
% САМОПРОВЕРКА
% ============================================================

\section*{5. Самопроверка перед ответом}
\sectionline

\begin{itemize}[leftmargin=8mm,itemsep=4pt]

\checkitem{Я разделил движение на отдельные участки.}

\checkitem{Я понимаю, что дано на каждом участке: путь, время, скорость, доля или процент.}

\checkitem{Я перевёл слова условия в выражения вида $S/2$, $2T/3$, $0{,}2S$.}

\checkitem{Я восстановил недостающие величины по формулам $S=vt$ или $t=S/v$.}

\checkitem{Я использовал формулу $v_{\text{ср}}=S_{\Sigma}/t_{\Sigma}$.}

\checkitem{Я аккуратно сложил дроби и проверил знаменатели.}

\checkitem{Я сокращал только общие множители.}

\checkitem{Я привёл единицы измерения к единому виду.}

\checkitem{Полученная средняя скорость лежит между минимальной и максимальной скоростями движения на участках.}

\checkitem{В ответе записаны число и единица измерения.}

\end{itemize}

% ============================================================
% БЛОК-СХЕМА
% ============================================================

\clearpage
\thispagestyle{empty}

\begin{center}
{\Large\bfseries\color{darkaccent}
Алгоритм решения задач на среднюю скорость
}
\end{center}

\vspace{6mm}

\begin{center}
\begin{tikzpicture}[node distance=8mm and 16mm]

\node[flowstart] (start)
{Начало};

\node[flowproc,below=of start] (split)
{
Разбейте движение на участки и выпишите для каждого участка известные $S_i$, $t_i$, $v_i$, доли или проценты.
};

\node[flowproc,below=of split] (translate)
{
Переведите словесные данные в математические выражения: например, $S/2$, $2T/3$, $20\%\,S=0{,}2S$.
};

\node[flowdec,below=11mm of translate] (missing)
{
Есть неизвестные
$S_i$ или $t_i$?
};

\node[
  flowproc,
  right=11mm of missing,
  text width=46mm
] (restore)
{
Восстановите их:
$S_i=v_it_i$
или
$t_i=S_i/v_i$.
};

\node[flowproc,below=11mm of missing] (sum)
{
Найдите полный путь
$S_{\Sigma}=\sum S_i$
и полное время
$t_{\Sigma}=\sum t_i$.
};

\node[flowproc,below=of sum] (average)
{
Вычислите
$\displaystyle v_{\text{ср}}=\frac{S_{\Sigma}}{t_{\Sigma}}$.
};

\node[flowproc,below=of average] (check)
{
Проверьте знаменатели, единицы измерения и физический смысл ответа.
};

\node[flowstart,below=of check] (finish)
{Ответ};

\draw[flowline] (start) -- (split);
\draw[flowline] (split) -- (translate);
\draw[flowline] (translate) -- (missing);
\draw[flowline] (missing) -- node[left,font=\small]{нет} (sum);
\draw[flowline] (missing.east) -- node[above,font=\small]{да} (restore.west);
\draw[flowline] (restore.south) |- (sum.east);
\draw[flowline] (sum) -- (average);
\draw[flowline] (average) -- (check);
\draw[flowline] (check) -- (finish);

\end{tikzpicture}
\end{center}

\clearpage

% ============================================================
% ТРЕНИРОВОЧНЫЙ БЛОК
% ============================================================

\section*{6. Тренировочный блок --- Путь Б}
\sectionline

Для каждой задачи оформляйте «Дано», основную формулу, промежуточные вычисления и ответ с единицами измерения.

\subsection*{27.09.2026}

\begin{enumerate}[
  label=\textbf{\arabic*.},
  leftmargin=8mm,
  itemsep=8pt
]

\item
Первые $3$ ч автомобиль двигался со скоростью $60$ км/ч, следующие $2$ ч --- со скоростью $30$ км/ч. Найдите среднюю скорость автомобиля на всём пути.

\item
Первую половину пути велосипедист ехал со скоростью $6$ км/ч, вторую половину --- со скоростью $12$ км/ч. Найдите среднюю скорость велосипедиста.

\item
Первую половину всего времени тело двигалось со скоростью $24$ км/ч, вторую половину времени --- со скоростью $72$ км/ч. Найдите среднюю скорость тела.

\end{enumerate}

\subsection*{28.09.2026}

\begin{enumerate}[
  label=\textbf{\arabic*.},
  leftmargin=8mm,
  itemsep=8pt
]

\item
Первые $20\%$ пути тело двигалось со скоростью $10$ м/с, следующие $50\%$ --- со скоростью $20$ м/с, оставшиеся $30\%$ --- со скоростью $15$ м/с. Найдите среднюю скорость тела на всём пути.

\item
Первые $\dfrac{2}{3}$ всего времени автомобиль двигался со скоростью $30$ км/ч, оставшуюся треть времени --- со скоростью $90$ км/ч. Найдите среднюю скорость автомобиля.

\item
Первую четверть пути тело двигалось со скоростью $5$ м/с, оставшиеся $\dfrac{3}{4}$ пути --- со скоростью $15$ м/с. Найдите среднюю скорость тела.

\end{enumerate}

% ============================================================
% ОТВЕТЫ
% ============================================================

\clearpage

\section*{7. Ответы к тренировочному блоку}
\sectionline

\renewcommand{\arraystretch}{1.35}

\begin{table}[h]
\centering
\begin{tabularx}{\linewidth}{@{}l c X@{}}
\toprule
\textbf{Дата} & \textbf{№} & \textbf{Ответ} \\
\midrule
27.09.2026 & 1 & $48$ км/ч \\
27.09.2026 & 2 & $8$ км/ч \\
27.09.2026 & 3 & $48$ км/ч \\
\midrule
28.09.2026 & 1 & $\dfrac{200}{13}$ м/с $\approx15{,}38$ м/с \\
28.09.2026 & 2 & $50$ км/ч \\
28.09.2026 & 3 & $10$ м/с \\
\bottomrule
\end{tabularx}
\end{table}

\vfill

\begin{center}
\small\color{softgray}
Если средняя скорость получилась меньше минимальной или больше максимальной скорости движения на участках, вычисления следует перепроверить.
\end{center}

\end{document}